% ---- pg-057 （书页 45） ----

乘积 $\bm{b}'\bm{c}$ 是一个 $1\times 1$ 向量或一个数,在这里等于 $-13$.
\[
\bm{b}\bm{c}' =
\begin{bmatrix}
7 \\
-3 \\
6
\end{bmatrix}
\begin{bmatrix}
5 & 8 & -4
\end{bmatrix}
=
\begin{bmatrix}
35 & 56 & -28 \\
-15 & -24 & 12 \\
30 & 48 & -24
\end{bmatrix}
\]

乘积 $\bm{b}\bm{c}'$ 是一矩阵,其行数等于 $\bm{b}$ 的维数,列数等于 $\bm{c}$ 的维数.这个积与 $\bm{b}'\bm{c}$ 不同,$\bm{b}'\bm{c}$ 是一个数.

\[
\bm{d}'\bm{A}\bm{b} =
\begin{bmatrix}
2 & 9
\end{bmatrix}
\begin{bmatrix}
1 & -2 & 3 \\
2 & 4 & -1
\end{bmatrix}
\begin{bmatrix}
7 \\
-3 \\
6
\end{bmatrix}
=
\begin{bmatrix}
26
\end{bmatrix}
\]

乘积 $\bm{d}'\bm{A}\bm{b}$ 是一个 $1\times 1$ 向量或一个数,在这里等于 26.\hfill$\blacksquare$

方阵在统计方法的开发中将具有特殊的重要性.对一个方阵,如果 $\bm{A}=\bm{A}'$ 或对所有 $i$ 和 $j$,$a_{ij}=a_{ji}$,则称此方阵为对称阵.

\mvbold{例 2.7}\quad (对称阵)

矩阵
\[
\begin{bmatrix}
3 & 5 \\
5 & -2
\end{bmatrix}
\]
是对称阵; 矩阵
\[
\begin{bmatrix}
3 & 6 \\
4 & -2
\end{bmatrix}
\]
则不是对称阵.\hfill$\blacksquare$

当两个方阵 $\bm{A}$ 和 $\bm{B}$ 维数相同时,两个乘积 $\bm{A}\bm{B}$ 和 $\bm{B}\bm{A}$ 均有定义,尽管它们未必相等(见补充 2A).若以 $\bm{I}$ 表示主对角线上元素均为 1,其余元素均为 0 的方阵,则按照矩阵乘法的定义,$\bm{A}\bm{I}$ 的第 $(i,j)$ 个元素为 $a_{i1}\times 0+\cdots+a_{i,i-1}\times 0+a_{ii}\times 1+a_{i,i+1}\times 0+\cdots+a_{ik}\times 0=a_{ij}$,于是 $\bm{A}\bm{I}=\bm{A}$.类似有 $\bm{I}\bm{A}=\bm{A}$.因此,对任何 $\bm{A}$ 有
\phantomsection\label{eq:2-11}
\begin{equation}
\underset{(k\times k)}{\bm{I}}\,\underset{(k\times k)}{\bm{A}} = \underset{(k\times k)}{\bm{A}}\,\underset{(k\times k)}{\bm{I}} = \underset{(k\times k)}{\bm{A}} \tag{2-11}
\end{equation}

矩阵 $\bm{I}$ 就像数字 1 在普通乘法中所起的作用一样 $(1\cdot a=a\cdot 1=a)$,因此称为单位矩阵.

如果 $a\neq 0$,就存在一个倒数 $a^{-1}$,能使 $a^{-1}a=aa^{-1}=1$,这个基本的数量关系在矩阵代数中有如下推广: 如果存在矩阵 $\bm{B}$,使得
\[
\underset{(k\times k)}{\bm{B}}\,\underset{(k\times k)}{\bm{A}} = \underset{(k\times k)}{\bm{A}}\,\underset{(k\times k)}{\bm{B}} = \underset{(k\times k)}{\bm{I}}
\]
则称 $\bm{B}$ 为 $\bm{A}$ 的逆矩阵,并记作 $\bm{A}^{-1}$.

存在逆矩阵的条件是: $\bm{A}$ 的 $k$ 个列 $\bm{a}_1,\bm{a}_2,\cdots,\bm{a}_k$ 线性无关.就是说,$\bm{A}^{-1}$ 的存在等价于
\phantomsection\label{eq:2-12}
\begin{equation}
c_1\bm{a}_1 + c_2\bm{a}_2 + \cdots + c_k\bm{a}_k = \bm{0} \quad\text{仅当}\quad c_1 = \cdots = c_k = 0 \text{ 时} \tag{2-12}
\end{equation}
(见补充 2A 中结论 2A.9.)

\mvbold{例 2.8}\quad (逆矩阵的存在性)

对于
\[
\underset{(2\times 2)}{\bm{A}} =
\begin{bmatrix}
3 & 2 \\
4 & 1
\end{bmatrix}
\]
可以验证
\[
\begin{bmatrix}
-0.2 & 0.4 \\
0.8 & -0.6
\end{bmatrix}
\begin{bmatrix}
3 & 2 \\
4 & 1
\end{bmatrix}
=
\begin{bmatrix}
(-0.2)\times 3 + \quad 0.4\times 4 & (-2)\times 2 + \quad 0.4\times 1 \\
0.8\times 3+(-0.6)\times 4 & 0.8\times 2+(-0.6)\times 1
\end{bmatrix}
\]
